% =============================================================================
% propwrap — Technical Report
% Industry-style engineering documentation (ideal CEA thermochemistry library)
% Compile: pdflatex propwrap_technical_report.tex  (run twice for ToC/refs)
% =============================================================================
\documentclass[11pt,a4paper]{report}

% --- Encoding & fonts ---
\usepackage[T1]{fontenc}
\usepackage[utf8]{inputenc}
\usepackage{lmodern}
\usepackage{microtype}

% --- Page & layout ---
\usepackage[margin=25mm,headheight=14pt]{geometry}
\usepackage{setspace}
\onehalfspacing
\usepackage{parskip}
\usepackage{booktabs}
\usepackage{longtable}
\usepackage{array}
\usepackage{tabularx}
\usepackage{multirow}
\usepackage{enumitem}
\usepackage{fancyhdr}
\usepackage{titlesec}
\usepackage{tocloft}
\usepackage{caption}
\usepackage{subcaption}
\usepackage{float}
\usepackage{graphicx}
\usepackage{xcolor}
\usepackage{hyperref}
\usepackage{amsmath,amssymb,amsfonts,bm}
\usepackage{siunitx}
\usepackage{listings}
\usepackage{csquotes}
\usepackage[numbers,sort&compress]{natbib}
\usepackage{appendix}

\sisetup{
  per-mode = symbol,
  detect-all,
  inter-unit-product = \ensuremath{{}\cdot{}}
}

\hypersetup{
  colorlinks=true,
  linkcolor=black!70!blue,
  citecolor=black!60!green,
  urlcolor=black!50!blue,
  pdftitle={propwrap Technical Report},
  pdfauthor={Shanthosh K V},
  pdfsubject={Liquid bipropellant thermochemistry software}
}

\pagestyle{fancy}
\fancyhf{}
\fancyhead[L]{\small\textit{propwrap} Technical Report}
\fancyhead[R]{\small\thepage}
\fancyfoot[C]{\footnotesize Independent open-source project --- not flight-certified}
\renewcommand{\headrulewidth}{0.4pt}

\titleformat{\chapter}[display]
  {\normalfont\huge\bfseries}{\chaptertitlename\ \thechapter}{16pt}{\Huge}
\titlespacing*{\chapter}{0pt}{-10pt}{24pt}

% --- Listings style (code) ---
\lstdefinestyle{py}{
  language=Python,
  basicstyle=\ttfamily\footnotesize,
  keywordstyle=\bfseries,
  commentstyle=\itshape\color{black!50},
  stringstyle=\color{black!40!blue},
  showstringspaces=false,
  breaklines=true,
  frame=single,
  rulecolor=\color{black!20},
  backgroundcolor=\color{black!3},
  tabsize=4,
  morekeywords={True,False,None}
}
\lstset{style=py}

% --- Custom commands ---
\newcommand{\pw}{\texttt{propwrap}}
\newcommand{\cea}{\textsc{cea}}
\newcommand{\vect}[1]{\boldsymbol{#1}}
\newcommand{\of}{\mathrm{O/F}}
\newcommand{\Isp}{I_{\mathrm{sp}}}
\newcommand{\cstar}{c^{*}}
\newcommand{\Cf}{C_{f}}
\newcommand{\ve}{v_{e}}
\newcommand{\Pc}{P_{c}}
\newcommand{\Pe}{P_{e}}
\newcommand{\Pa}{P_{a}}
\newcommand{\eps}{\varepsilon}
\newcommand{\gzero}{g_{0}}

\begin{document}

% =============================================================================
% TITLE PAGE
% =============================================================================
\begin{titlepage}
  \centering
  \vspace*{1.2cm}
  {\Large\bfseries Technical Report}\\[0.6cm]
  {\huge\bfseries propwrap}\\[0.35cm]
  {\Large A Software Library for Liquid Bipropellant\\
  Thermochemical Performance Evaluation}\\[1.2cm]

  \begin{tabular}{rl}
    \textbf{Document type} & Engineering software technical report \\
    \textbf{Software version} & 0.1.3 \\
    \textbf{Distribution} & PyPI package \texttt{propwrap}; MIT source license \\
    \textbf{Author} & Shanthosh K.\ V.\ (\texttt{shanthoshkv@gmail.com}) \\
    \textbf{Repository} & \url{https://github.com/shanthoshkv/propwrap} \\
    \textbf{Project website} & \url{https://aboutkvs.vercel.app/propwrap.html} \\
    \textbf{Date} & 30 July 2026 \\
  \end{tabular}

  \vfill
  \begin{minipage}{0.92\textwidth}
    \small
    \textbf{Disclaimer.}
    This document describes theoretical, one-dimensional chemical equilibrium
    performance as implemented by NASA~\cea{} (via RocketCEA). Numbers produced
    by the software are \emph{not} flight-delivered engine performance and do not
    constitute certification evidence. The project is independent open-source work
    and is \textbf{not affiliated with} ISRO, DRDO, NAL, HAL, NASA, or any other
    agency or company.
  \end{minipage}

  \vspace{1.2cm}
  {\footnotesize MIT License for \pw{} source code.
  Third-party engines (RocketCEA/GPL-family, Cantera, etc.) retain their own licenses.}
\end{titlepage}

% =============================================================================
% FRONT MATTER
% =============================================================================
\pagenumbering{roman}
\tableofcontents
\newpage
\listoffigures
\listoftables
\newpage

\chapter*{Abstract}
\addcontentsline{toc}{chapter}{Abstract}
\noindent
Liquid rocket design begins with propellant thermochemistry: mixture ratio,
chamber pressure, and nozzle expansion set the ideal specific impulse,
characteristic velocity, product temperature, and molecular weight that every
downstream subsystem inherits. NASA's Chemical Equilibrium with Applications
(\cea{}) program remains the standard way to compute those numbers, but
day-to-day use is still awkward---English units leak into SI notebooks, scripts
drift, and multi-propellant trades are often run at a single forced mixture
ratio that is fair to none of the pairs.

\pw{} is a Python~3.10+ library that wraps NASA~\cea{} through RocketCEA and
exposes a typed, SI-first public API. It computes vacuum and sea-level
specific impulse (shifting and frozen chemistry), effective exhaust velocity,
characteristic velocity, thrust coefficients, station temperatures and
pressures, product $\gamma$, bulk liquid density, and density impulse. Higher-level
workflows support O/F sweeps, fair multi-pair trades (each pair at its own
optimum O/F), multi-component blends, unit conversion, SQLite result caching,
a command-line interface, optional Cantera frozen-$\gamma$ cross-checks, and
student-oriented lab report packs.

This report documents the engineering motivation, governing performance
relations, software architecture, unit system, public interfaces, validation
strategy (including published RocketCEA goldens and physics identity tests),
worked examples, and known limitations. The intent is that a propulsion engineer
or graduate student can read this document, install the package from PyPI, and
reproduce the same theoretical numbers with a clear understanding of what they
do---and do not---mean.

\vspace{0.8em}
\noindent\textbf{Keywords:}
liquid rocket propellants; NASA CEA; RocketCEA; specific impulse; characteristic
velocity; density impulse; SI units; software validation; Python.

\newpage
\chapter*{Nomenclature}
\addcontentsline{toc}{chapter}{Nomenclature}

\section*{Symbols}
\begin{longtable}{@{}p{2.4cm}p{11.5cm}@{}}
$A_e$, $A_t$ & Exit and throat areas [\si{\metre\squared}] \\
$\Cf$ & Thrust coefficient [---] \\
$\cstar$ & Characteristic velocity [\si{\metre\per\second}] \\
$c_p$, $c_v$ & Specific heats at constant pressure / volume \\
$F$ & Thrust [\si{\newton}] \\
$\gzero$ & Standard gravity, \SI{9.80665}{\metre\per\second\squared} \\
$\gamma$ & Isentropic exponent (performance $\gamma$ from \cea{}) \\
$\Isp$ & Specific impulse [\si{\second}] \\
$\dot{m}$ & Mass flow rate [\si{\kilogram\per\second}] \\
$M_w$ & Molecular weight [\si{\kilogram\per\kilo\mole}] \\
$\of$ & Mixture ratio, $m_{\mathrm{ox}}/m_{\mathrm{fuel}}$ [---] \\
$\Pc$, $\Pe$, $\Pa$ & Chamber, exit, ambient pressure [\si{\pascal}] \\
$R$ & Specific gas constant; $R_u$ universal gas constant \\
$T_c$, $T_t$, $T_e$ & Chamber, throat, exit temperatures [\si{\kelvin}] \\
$\ve$ & Effective exhaust velocity [\si{\metre\per\second}] \\
$\eps$ & Nozzle area ratio $A_e/A_t$ [---] \\
$\eta_{c^{*}}$, $\eta_{\Cf}$ & Optional efficiency factors applied to $\cstar$ and $\Cf$ \\
$\rho$ & Density [\si{\kilogram\per\metre\cubed}] \\
$\mu$, $k$, $\mathrm{Pr}$ & Viscosity, thermal conductivity, Prandtl number \\
\end{longtable}

\section*{Abbreviations}
\begin{longtable}{@{}p{2.8cm}p{11.1cm}@{}}
\cea{} & Chemical Equilibrium with Applications (NASA) \\
CLI & Command-line interface \\
DfAM & Design for additive manufacturing (context only; not in \pw{}) \\
LOX & Liquid oxygen \\
LH$_2$ & Liquid hydrogen \\
MMH & Monomethylhydrazine \\
N$_2$O$_4$ & Nitrogen tetroxide \\
RP-1 & Rocket Propellant-1 (kerosene-class) \\
SI & International System of Units \\
UDMH & Unsymmetrical dimethylhydrazine \\
\end{longtable}

\newpage
\pagenumbering{arabic}

% =============================================================================
\chapter{Introduction}
% =============================================================================

\section{Why propellant numbers matter early}

Before anyone sizes a turbopump or draws a cooling channel, someone has to
answer a deceptively simple question: \emph{if this fuel and oxidizer burn at
this mixture ratio and chamber pressure, and expand through this area ratio,
what ideal performance do we get?}

Those answers---$\Isp$, $\cstar$, $T_c$, $\gamma$, $M_w$---propagate into every later
calculation. A few seconds of $\Isp$ error looks small on a slide and expensive
on a mass budget. A wrong $\gamma$ or molecular weight spoils a nozzle contour
study before the first CAD feature is drawn. In practice the industry has used
NASA~\cea{} for decades precisely because the equilibrium chemistry problem is
too messy to reinvent casually~\citep{gordon1994cea,gordon1996cea}.

\section{The practical problem \pw{} addresses}

\cea{} itself is solid. The surrounding workflow is often not:

\begin{itemize}[leftmargin=*]
  \item Scripts and notebook cells accumulate English units (psia, Rankine,
        ft/s) next to SI formulas.
  \item Results live as parsed strings or ad-hoc dictionaries.
  \item ``Fair'' propellant trades sometimes force one O/F onto hydrogen and
        kerosene alike---chemically unfair and easy to publish by mistake.
  \item Re-running last semester's case means hoping the same CEA deck still
        exists on someone's laptop.
\end{itemize}

\pw{} exists to make the propellant layer boring and reliable: one importable
library, SI public results, typed models, reproducible tests, and workflows that
make the right comparison the easy comparison.

\section{What this report covers}

Chapters~\ref{ch:theory}--\ref{ch:units} state the performance theory and the
unit conventions the software enforces. Chapter~\ref{ch:arch} describes
architecture and major modules. Chapter~\ref{ch:api} summarises the public API
and CLI. Chapter~\ref{ch:validation} documents validation claims and automated
tests. Chapter~\ref{ch:examples} walks through representative cases.
Chapter~\ref{ch:limits} states limitations honestly. References and appendices
collect sources and package layout.

\section{Project identity}

\begin{itemize}[leftmargin=*]
  \item \textbf{Name:} \pw{} (PyPI: \texttt{propwrap})
  \item \textbf{Language / requirement:} Python~$\geq 3.10$
  \item \textbf{Core engine:} NASA~\cea{} via RocketCEA~\citep{rocketcea}
  \item \textbf{Optional cross-check:} Cantera frozen $\gamma$~\citep{cantera}
  \item \textbf{Source license:} MIT for \pw{} code; RocketCEA is GPL-family---both
        must be reviewed for proprietary redistribution
  \item \textbf{Affiliation:} Independent open source; not an agency product
\end{itemize}

% =============================================================================
\chapter{Scope of the software}
\label{ch:scope}
% =============================================================================

\section{In scope}

\pw{} addresses \textbf{liquid bipropellant thermochemical performance} under
ideal one-dimensional \cea{} assumptions:

\begin{itemize}[leftmargin=*]
  \item Shifting (equilibrium) and frozen composition rocket performance
  \item Vacuum and sea-level specific impulse
  \item Effective exhaust velocity $\ve=\Isp\gzero$
  \item Characteristic velocity $\cstar$ and thrust coefficient $\Cf$
  \item Chamber / throat / exit temperatures and pressures
  \item Product $\gamma$, molecular weight, selected transport properties
  \item Liquid bulk density and density impulse $\rho\Isp$
  \item O/F, $\Pc$, and $\eps$ parametric sweeps
  \item Multi-pair trades at each pair's own optimum O/F
  \item Multi-component weight-percent blends and a propellant registry
  \item SI converters, CLI, SQLite cache, student homework packs
\end{itemize}

\section{Out of scope}

The library is deliberately \emph{not} a full engine design code:

\begin{itemize}[leftmargin=*]
  \item No injector CFD, droplet combustion, or finite-rate kinetics
  \item No regenerative cooling channel sizing or structural margins
  \item No turbopump or feed-system geometry
  \item No trajectory, vehicle $\Delta v$, or stage optimisation
  \item No claim that theoretical $\Isp$ equals flight-delivered $\Isp$
  \item No flight certification or agency endorsement
\end{itemize}

Those tools can sit \emph{on top} of \pw{} and pull $\gamma$, $\cstar$, $T$, and
$\Isp$ from one tested place.

\section{Truthfulness of numbers}

Unless the user applies optional efficiency knobs $\eta_{c^{*}}$ and
$\eta_{\Cf}$, every performance number is \textbf{ideal theoretical \cea{}}.
Real engines deliver less. That distinction is written into the documentation,
lab packs, and this report on purpose.

% =============================================================================
\chapter{Theoretical background}
\label{ch:theory}
% =============================================================================

This chapter recalls the classical rocket performance definitions that \cea{}
implements and that \pw{} exposes. Derivations follow standard propulsion
texts~\citep{sutton2017,huzel1992} and the NASA~\cea{} rocket problem
statement~\citep{gordon1994cea}.

\section{Thrust and effective exhaust velocity}

The axial thrust of an ideal rocket in ambient pressure $\Pa$ is
\begin{equation}
  F = \dot{m}\,\ve' + (\Pe - \Pa)\,A_e,
  \label{eq:thrust}
\end{equation}
where $\ve'$ is the exit velocity relative to the vehicle. Engineers often fold
pressure thrust into an \emph{effective} exhaust velocity $\ve$ such that
\begin{equation}
  F = \dot{m}\,\ve.
  \label{eq:F_mdot_ve}
\end{equation}
In vacuum ($\Pa=0$) with the usual definitions used in equilibrium performance
codes, vacuum specific impulse and effective exhaust velocity are related by
the conventional standard gravity:
\begin{equation}
  \boxed{\ve = \Isp\,\gzero, \qquad \gzero = \SI{9.80665}{\metre\per\second\squared}.}
  \label{eq:ve_isp}
\end{equation}
Equation~\eqref{eq:ve_isp} is not a local gravitational acceleration on a
particular launch pad; it is the metrological convention that converts
$\Isp$ in seconds to velocity in SI~\citep{bipm_g0}. \pw{} uses this constant
exactly and tests that $\ve$ fields match $\Isp\times\gzero$ bit-for-bit within
floating-point expectations.

\section{Specific impulse}

Specific impulse is thrust per unit weight-flow using $\gzero$:
\begin{equation}
  \Isp = \frac{F}{\dot{m}\,\gzero} = \frac{\ve}{\gzero}.
  \label{eq:isp_def}
\end{equation}
\cea{} reports shifting (equilibrium composition during expansion) and frozen
(composition fixed at the throat or chamber, depending on the option) values.
For ideal cases of practical interest one expects
\begin{equation}
  \Isp^{\mathrm{shifting}} \ge \Isp^{\mathrm{frozen}},
  \qquad
  \Isp^{\mathrm{vac}} \ge \Isp^{\mathrm{SL}}
  \label{eq:isp_order}
\end{equation}
for the same $\Pc$ and $\eps$. \pw{} asserts these inequalities on locked
regression cases.

\section{Characteristic velocity}

Characteristic velocity measures chamber energy release and throat delivery,
independent of the nozzle expansion:
\begin{equation}
  \cstar = \frac{\Pc\,A_t}{\dot{m}}.
  \label{eq:cstar}
\end{equation}
In SI, kerolox-class theoretical $\cstar$ is typically of order
\SIrange{1700}{1900}{\metre\per\second}, not the numerically similar English
value in ft/s. A recurring software bug class is leaving $\cstar$ in ft/s while
labelling it m/s; \pw{} converts at the \cea{} adapter and tests SI bands.

\section{Thrust coefficient}

The thrust coefficient links chamber pressure and throat area to thrust:
\begin{equation}
  \Cf = \frac{F}{\Pc\,A_t}.
  \label{eq:cf_def}
\end{equation}
Combining \eqref{eq:F_mdot_ve}, \eqref{eq:isp_def}, and \eqref{eq:cstar} yields
the vacuum identity used throughout the code base:
\begin{equation}
  \boxed{\Cf^{\mathrm{vac}} = \frac{\ve^{\mathrm{vac}}}{\cstar} = \frac{\Isp^{\mathrm{vac}}\,\gzero}{\cstar}.}
  \label{eq:cf_ve_cstar}
\end{equation}
\pw{} recomputes and checks \eqref{eq:cf_ve_cstar} both in unit tests and in an
optional runtime sanity check.

\section{Nozzle area ratio and pressure ratio}

The geometric expansion ratio is
\begin{equation}
  \eps = \frac{A_e}{A_t}.
  \label{eq:eps}
\end{equation}
For isentropic flow of a calorically perfect gas one can write closed-form
Mach--area and pressure ratios in terms of $\gamma$. Real \cea{} products are
multi-species and temperature-dependent; the code solves equilibrium (or frozen)
chemistry rather than assuming constant $\gamma$. Still, the perfect-gas
relations remain useful pedagogy and sanity bounds:
\begin{align}
  \frac{A_e}{A_t}
  &= \frac{1}{M_e}\left[
      \frac{1+\dfrac{\gamma-1}{2}M_e^{2}}{1+\dfrac{\gamma-1}{2}}
     \right]^{\frac{\gamma+1}{2(\gamma-1)}},
  \label{eq:area_mach}\\[0.5em]
  \frac{\Pc}{\Pe}
  &= \left(1+\frac{\gamma-1}{2}M_e^{2}\right)^{\frac{\gamma}{\gamma-1}}.
  \label{eq:pcpe_perfect}
\end{align}
\pw{} stores $\Pc$ and $\Pe$ in pascals and verifies that $\Pc/\Pe$ matches the
ratio reported by the \cea{} path for the same case.

\section{Ideal-gas density check}

At the chamber, a useful consistency check on molecular weight and temperature
is
\begin{equation}
  \rho_c \approx \frac{\Pc\,M_w}{R_u\,T_c},
  \label{eq:rho_ig}
\end{equation}
with $R_u=\SI{8314.462618}{\joule\per\kilo\mole\per\kelvin}$ as coded in
\pw{}~\citep{codata}. Residual differences can appear when transport or
multi-phase bookkeeping differs; the automated test tolerates a small relative
error on a standard LOX/RP-1 point.

\section{Mixture ratio}

The oxidizer-to-fuel mixture ratio on a mass basis is
\begin{equation}
  \of = \frac{\dot{m}_{\mathrm{ox}}}{\dot{m}_{\mathrm{fuel}}}.
  \label{eq:of}
\end{equation}
Stoichiometric $\of$ depends on chemistry; maximum vacuum $\Isp$ is usually
fuel-rich of stoich for carbon-containing fuels. \pw{} can mark stoich on O/F
plots and can optimise $\of$ numerically for a fixed $\Pc$ and $\eps$.

\section{Liquid bulk density and density impulse}

For tankage- and volume-limited stages, density-weighted performance matters.
With liquid densities $\rho_f$, $\rho_o$ and mass mixture ratio $\of$,
\begin{equation}
  \rho_{\mathrm{bulk}}
  = \frac{(\of+1)\,\rho_f\,\rho_o}{\of\,\rho_f + \rho_o}
  \label{eq:rho_bulk}
\end{equation}
(standard volume-additive blend of two incompressible liquids). Density impulse
is then
\begin{equation}
  I_{\rho} = \rho_{\mathrm{bulk}}\,\Isp.
  \label{eq:density_isp}
\end{equation}
The $\of$ that maximises $\Isp$ is not always the $\of$ that maximises
$I_{\rho}$. \pw{} can report both curves and rank propellant pairs by either
metric.

\section{Shifting versus frozen expansion}

\begin{description}[style=nextline,leftmargin=1.2cm]
  \item[Shifting (equilibrium)] Composition re-equilibrates as $T$ and $P$ drop
        in the nozzle. Dissociation energy can recombine, often raising $\Isp$.
  \item[Frozen] Composition fixed after a defined station (commonly throat).
        More conservative for some systems; lower $\Isp$ when recombination would
        have helped.
\end{description}
Neither model is ``the flight answer.'' Flight includes kinetics, boundary
layers, divergence, finite-rate chemistry, and mixture maldistribution.
Shifting and frozen are brackets on the chemistry story, not a substitute for
hot-fire data~\citep{sutton2017}.

\section{Optional efficiency factors}

For rough knock-downs toward delivered performance, some workflows multiply
\begin{equation}
  \cstar_{\mathrm{eff}} = \eta_{c^{*}}\,\cstar,
  \qquad
  C_{f,\mathrm{eff}} = \eta_{C_f}\,C_f,
  \label{eq:eta}
\end{equation}
and recombine into an estimated delivered $\Isp$. In \pw{} these factors are
explicit knobs, not calibrated engine models. They must not be confused with a
validated power-balance simulation.

\section{Note on $\gamma$ versus $c_p/c_v$}

\cea{} publishes a performance $\gamma$ used for nozzle relations. Chamber
transport $c_p$ may include equilibrium dissociation contributions. Therefore
\begin{equation}
  \gamma \;\neq\; \frac{c_p}{c_p - R}
  \label{eq:gamma_not_cp}
\end{equation}
exactly is \textbf{expected chemistry}, not evidence of a unit-conversion bug.
\pw{} documents this distinction so users do not ``fix'' a non-problem.

% =============================================================================
\chapter{Units and metrology}
\label{ch:units}
% =============================================================================

\section{Public SI contract}

Outside the \cea{} adapter, \pw{} exposes:

\begin{center}
\begin{tabular}{@{}lll@{}}
\toprule
Quantity & Unit & Notes \\
\midrule
Pressure & \si{\pascal} & also \texttt{pc\_bar} view / input helpers \\
Temperature & \si{\kelvin} & \\
Velocity, $\cstar$, $\ve$ & \si{\metre\per\second} & \\
Density & \si{\kilogram\per\metre\cubed} & \\
$\Isp$ & \si{\second} & rocketry convention \\
Density impulse & \si{\second\kilogram\per\metre\cubed} & $\rho\Isp$ \\
$c_p$ & \si{\joule\per\kilogram\per\kelvin} & \\
Viscosity & \si{\pascal\second} & \\
\bottomrule
\end{tabular}
\end{center}

RocketCEA continues to speak English units internally (psia, Rankine, ft/s,
etc.). Conversion happens once in \texttt{cea\_backend}; public models store SI.

\section{Pressure inputs}

Users may enter chamber pressure as
\begin{align}
  \Pc &= \texttt{pc} = \texttt{pc\_pa} && \text{[\si{\pascal}]},\\
  \Pc &= \texttt{pc\_bar}\times 10^{5} && \text{[bar exact]},\\
  \Pc &= \texttt{pc\_mpa}\times 10^{6} && \text{[MPa]}.
\end{align}
The bar is exact by SI definition:
\begin{equation}
  1\,\mathrm{bar} = 10^{5}\,\mathrm{Pa}.
  \label{eq:bar}
\end{equation}
A common student error is \texttt{pc=70} intending 70~bar. That is 70~Pa---near
vacuum---and produces nonsense. \pw{} rejects pathologically low $\Pc$ values
that almost always mean ``forgot the unit.''

\section{Constants coded in the library}

\begin{center}
\begin{tabular}{@{}lll@{}}
\toprule
Symbol & Value & Role \\
\midrule
$\gzero$ & \SI{9.80665}{\metre\per\second\squared} & $\ve=\Isp\gzero$ \\
$R_u$ & \SI{8314.462618}{\joule\per\kilo\mole\per\kelvin} & ideal-gas $\rho$ check \\
$1\,\mathrm{bar}$ & \SI{1e5}{\pascal} & \texttt{pc\_bar} \\
$1\,\mathrm{atm}$ & \SI{101325}{\pascal} & sea-level ambient reference \\
\bottomrule
\end{tabular}
\end{center}

\section{English $\leftrightarrow$ SI (adapter only)}

Representative conversions applied when reading RocketCEA outputs include
\begin{align}
  T[\mathrm{K}] &= T[\mathrm{R}]\times\frac{5}{9},\\
  P[\mathrm{Pa}] &= P[\mathrm{psia}]\times 6894.757293168,\\
  v[\mathrm{m/s}] &= v[\mathrm{ft/s}]\times 0.3048.
\end{align}
Users should not need these at the public API; they are listed for auditability.

% =============================================================================
\chapter{Software architecture}
\label{ch:arch}
% =============================================================================

\section{Design goals}

\begin{enumerate}[leftmargin=*]
  \item \textbf{Library, not ritual} --- importable, version-pinned, testable.
  \item \textbf{SI wall} --- English units stop at the \cea{} backend.
  \item \textbf{Typed results} --- Pydantic models; readable \texttt{print(result)}.
  \item \textbf{Fair trades by default} --- optimise O/F per pair before ranking.
  \item \textbf{Fail loudly} --- physics identities and pressure guards.
\end{enumerate}

\section{Layered structure}

Figure~\ref{fig:arch} sketches the call path.

\begin{figure}[H]
\centering
\begin{verbatim}
  User code / CLI / sibling tool
              |
              v
     Mixture / workflows / units
              |
              v
        Pydantic result models
              |
        +-----+------+
        |            |
        v            v
  cea_backend   cantera_backend (frozen gamma)
  (RocketCEA)         |
        |             |
        +------+------+
               |
               v
        SQLite cache (~/.propwrap/)
\end{verbatim}
\caption{Logical architecture of \pw{}. Public code never re-implements
equilibrium chemistry; it orchestrates RocketCEA/Cantera and polices units.}
\label{fig:arch}
\end{figure}

\section{Major modules}

\begin{center}
\begin{tabular}{@{}p{3.6cm}p{10.2cm}@{}}
\toprule
Module & Responsibility \\
\midrule
\texttt{propellant.py} & \texttt{Mixture} primary API; evaluate, sweeps, studies \\
\texttt{cea\_backend.py} & RocketCEA calls; English$\to$SI conversion \\
\texttt{cantera\_backend.py} & Frozen $\gamma$ from mapped major species \\
\texttt{models.py} & Pydantic result types (\texttt{PerformanceResult}, \ldots) \\
\texttt{units.py} & Constants and converters \\
\texttt{workflows.py} & \texttt{characterize}, \texttt{compare\_propellants}, blends \\
\texttt{trades.py} & Optimum-O/F trades; density-Isp curves \\
\texttt{registry.py} & Names, aliases, densities, cryo defaults \\
\texttt{blends.py} & Multi-component weight-percent cards \\
\texttt{cache.py} & SQLite result cache \\
\texttt{cli.py} / \texttt{\_\_main\_\_.py} & Console entry points \\
\texttt{reports.py} & Markdown/CSV lab packs \\
\texttt{plotting.py} & Optional dark-theme figures (matplotlib) \\
\texttt{sanity.py} & Runtime physics identity checks \\
\bottomrule
\end{tabular}
\end{center}

\section{Dependencies}

Hard dependencies of version~0.1.x:

\begin{itemize}[leftmargin=*]
  \item \textbf{rocketcea} --- NASA~\cea{} FORTRAN via Python~\citep{rocketcea}
  \item \textbf{cantera} --- optional-use cross-check path; required install~\citep{cantera}
  \item \textbf{pydantic} $\geq 2$ --- result models
  \item \textbf{numpy}, \textbf{matplotlib} --- numerics and plotting
\end{itemize}

\section{Caching}

Repeated identical cases hit a local SQLite database under
\texttt{\textasciitilde/.propwrap/}. Cache keys include propellant identities,
$\of$, $\Pc$, $\eps$, and relevant options. Cache invalidation is available via
CLI and API. Caching never changes physics; it only avoids re-entering \cea{}.

% =============================================================================
\chapter{Public interface and workflows}
\label{ch:api}
% =============================================================================

\section{Installation}

From PyPI (recommended):
\begin{lstlisting}
python -m venv .venv
# Windows: .venv\Scripts\activate
pip install propwrap
propwrap run RP-1 LOX --of 2.56 --pc-bar 70 --eps 20
\end{lstlisting}

Requires Python~3.10+. On multi-Python Windows machines, prefer
\texttt{py -3.12 -m pip install propwrap} so an old 3.9 \texttt{pip} is not used.

\section{Minimal Python usage}

\begin{lstlisting}
from propwrap import Mixture, compare_propellants, convert

m = Mixture("RP-1", "LOX")
r = m.evaluate(of=2.56, pc_bar=70, eps=20)
print(r.isp_vac_shifting, r.ve_vac_shifting, r.c_star)

sw = m.scan_of((2.0, 3.2, 0.1), pc_bar=70, eps=20)
print(sw.optimum().of_ratio)

print(compare_propellants(
    ["RP-1/LOX", "CH4/LOX", "LH2/LOX"],
    pc_bar=70, eps=40,
))
print(convert(70, "bar", "Pa"))
\end{lstlisting}

Representative LOX/RP-1 point at $\of=2.56$, $\Pc=70\,\mathrm{bar}$,
$\eps=20$ (theoretical shifting vacuum):
\begin{equation}
  \Isp^{\mathrm{vac}} \approx \SI{343.7}{\second},\quad
  \ve \approx \SI{3370}{\metre\per\second},\quad
  \cstar \approx \SI{1798}{\metre\per\second}.
  \label{eq:kerolox_demo}
\end{equation}

\section{Primary API surface (stable for 0.1.x)}

\begin{itemize}[leftmargin=*]
  \item \texttt{Mixture} / \texttt{Propellant} --- evaluate, \texttt{scan\_of},
        density impulse, gamma profiles
  \item \texttt{compare\_propellants}, \texttt{characterize}, \texttt{define\_blend}
  \item \texttt{convert}, \texttt{propwrap.units}
  \item \texttt{PerformanceResult} fields documented in the API reference
  \item CLI: \texttt{run}, \texttt{scan-of}, \texttt{characterize},
        \texttt{compare-pairs}, \texttt{homework}, \ldots
\end{itemize}

May still change within 0.1.x: plot styling, Cantera cross-check details,
efficiency-knob presentation, internal cache layout.

\section{Fair multi-pair trades}

A trade that ranks LOX/LH$_2$ against LOX/RP-1 at the same $\of=2.5$ is not a
propellant trade; it is a self-inflicted handicap on one chemistry. The default
workflow is:

\begin{enumerate}[leftmargin=*]
  \item Fix $\Pc$ and $\eps$ (mission-class assumptions).
  \item For each pair, scan $\of$ and take the optimum (e.g.\ max vacuum
        shifting $\Isp$ or max $\rho\Isp$).
  \item Rank pairs on the chosen metric.
\end{enumerate}

That is what \texttt{compare\_propellants} encodes.

\section{Command-line interface}

\begin{lstlisting}[language=bash]
propwrap run RP-1 LOX --of 2.56 --pc-bar 70 --eps 20
propwrap scan-of RP-1 LOX --pc-bar 70 --eps 20 --range 2.0 3.5 0.1
propwrap compare-pairs --combos "RP-1/LOX,LH2/LOX,CH4/LOX" --pc-bar 70 --eps 40
propwrap homework kerolox --name YourName
\end{lstlisting}
Module form (0.1.1+): \texttt{python -m propwrap \ldots}

\section{Student tooling}

Coursework environments rarely match research desktops. \pw{} ships:

\begin{itemize}[leftmargin=*]
  \item Install notes for lab PCs
  \item Beginner walkthrough and concept notes (units, frozen/shifting, fair trades)
  \item One-page cheat sheet and sample lab assignment
  \item \texttt{propwrap homework} packs: summary, assumptions, CSV, figures
  \item Case presets: student lab / booster / upper stage
\end{itemize}

% =============================================================================
\chapter{Validation}
\label{ch:validation}
% =============================================================================

\section{What is claimed}

\begin{center}
\begin{tabular}{@{}p{7.2cm}cc@{}}
\toprule
Claim & Supported? & Mechanism \\
\midrule
Match RocketCEA/\cea{} for same inputs & Yes & Bit-level $\Isp$ tests \\
SI public units & Yes & Guards + $\rho$ check \\
$\ve=\Isp\gzero$ & Yes & Identity tests \\
$\Cf^{\mathrm{vac}}=\ve/\cstar$ & Yes & Identity tests \\
$\Pc/\Pe$ consistency & Yes & Identity tests \\
Equal to flight Merlin / RS-25 $\Isp$ & No & Explicitly excluded \\
Cantera $\gamma$ identical to \cea{} & No & Approximate frozen compare \\
\bottomrule
\end{tabular}
\end{center}

\section{Primary sources}

\subsection{NASA CEA methodology}

Gordon and McBride, NASA RP-1311 Parts~I and~II~\citep{gordon1994cea,gordon1996cea},
define the equilibrium/frozen rocket performance problem. \pw{} does not
re-derive the chemistry solver; it wraps the industry-standard implementation
and owns API integrity.

\subsection{RocketCEA published goldens}

The RocketCEA QuickStart publishes LOX/LH$_2$ vacuum $\Isp$ at
$\Pc=\SI{100}{psia}$, $\eps=40$~\citep{rocketcea}. \pw{} matches
\texttt{CEA\_Obj.get\_Isp} to machine precision for O/F~$=2\ldots 8$ on the
locked stack (verified 2026-07-29).

\begin{table}[H]
\centering
\caption{LOX/LH$_2$ golden vacuum $\Isp$ (RocketCEA docs conditions).}
\label{tab:lox_lh2_gold}
\begin{tabular}{@{}rrr@{}}
\toprule
O/F & $\Isp^{\mathrm{vac}}$ [s] & Role \\
\midrule
default \texttt{get\_Isp} & 374.3036\ldots & install smoke \\
2 & 424.36 & fuel-rich \\
4 & 453.13 & \\
5 & 453.24 & near peak \\
6 & 448.19 & \\
8 & 424.70 & ox-rich \\
\bottomrule
\end{tabular}
\end{table}

\subsection{Metrology}

$\gzero$ and the exact bar follow BIPM/SI conventions~\citep{bipm_g0}.

\section{Physics identities (must not drift)}

Automated tests (\texttt{tests/test\_physics\_identities.py}) lock:

\begin{enumerate}[leftmargin=*]
  \item Equation~\eqref{eq:ve_isp}
  \item Equation~\eqref{eq:cf_ve_cstar}
  \item $\Pc/\Pe$ vs stored pressures
  \item Chamber density check~\eqref{eq:rho_ig}
  \item $T_e < T_t < T_c$
  \item Typical $\gamma_e \ge \gamma_c$
  \item Shifting $\ge$ frozen; vac $\ge$ SL on locked cases
  \item $\cstar$ SI magnitude band
  \item Density impulse product~\eqref{eq:density_isp}
\end{enumerate}

\subsection{Bug fixed during physics audit}

\texttt{gamma\_and\_temp\_at\_eps} previously returned exit temperature twice
instead of chamber and exit. The bug is fixed and covered by tests (audit date
2026-07-29).

\section{Multi-propellant regression matrix}

Representative theoretical anchors (CEA via RocketCEA; not flight):

\begin{center}
\footnotesize
\begin{tabular}{@{}llrrr@{}}
\toprule
Pair & O/F & $\Pc$ & $\eps$ & $\Isp^{\mathrm{vac}}$ (shift.) \\
\midrule
RP-1/LOX & 2.56 & 70~bar & 20 & $\sim$343.7~s \\
RP-1/LOX & 2.3 & $\sim$1000~psia & 40 & $\sim$352.3~s \\
LH$_2$/LOX & 5.0 & $\sim$1000~psia & 40 & $\sim$454.9~s \\
LH$_2$/LOX (RS-25-class) & 5.5 & 100~bar & 69 & $\sim$463.7~s \emph{theory} \\
CH$_4$/LOX & 3.2 & $\sim$1000~psia & 40 & $\sim$367.6~s \\
MMH/N$_2$O$_4$ & 2.0 & $\sim$1000~psia & 40 & $\sim$338.6~s \\
\bottomrule
\end{tabular}
\end{center}

\textbf{Flight contrast (not used as code targets):}
RS-25 delivered vacuum $\Isp$ is about \SI{452}{\second}; pure CEA at similar
class conditions sits higher. Merlin-class flight kerolox $\Isp$ is well below
pure theoretical CEA. Mixing the two without labelling is a documentation
failure; the catalog marks the gap deliberately.

\section{Test suite}

On the development stack, the full suite reports on the order of \textbf{132}
passing tests spanning physics identities, validation catalog, known values,
units, cache, sweeps, registry/trades, student helpers, and UX guards.

\begin{lstlisting}[language=bash]
pytest tests/test_physics_identities.py tests/test_validation_suite.py -v
pytest -q
\end{lstlisting}

% =============================================================================
\chapter{Worked examples}
\label{ch:examples}
% =============================================================================

\section{Single-point kerolox}

\textbf{Inputs:} RP-1 / LOX, $\of=2.56$, $\Pc=70\,\mathrm{bar}$, $\eps=20$.

\textbf{Outputs (theoretical shifting, representative run):}
$\Isp^{\mathrm{vac}}\approx\SI{343.7}{\second}$,
$\ve\approx\SI{3370}{\metre\per\second}$,
$\cstar\approx\SI{1798}{\metre\per\second}$,
$T_c\sim\SI{3670}{\kelvin}$ class,
$\Pc$ stored as \SI{7e6}{\pascal}.

\textbf{Use:} classroom baseline; regression lock REG-RP1-LOX-70bar-eps20.

\section{O/F sweep}

Scan $\of\in[2.0,3.2]$ at fixed $\Pc$, $\eps$; locate maximum vacuum shifting
$\Isp$ and optionally maximum $\rho\Isp$. Expect the $\Isp$ peak fuel-rich of
stoichiometric for kerolox-class chemistry.

\section{Fair three-pair trade}

Compare RP-1/LOX, CH$_4$/LOX, and LH$_2$/LOX at the same $\Pc$ and $\eps$, each
at its own optimum $\of$. LH$_2$ leads on $\Isp$; dense hydrocarbons lead on
$\rho\Isp$. Both rankings are valid---they answer different vehicle questions.

\section{Homework pack}

\begin{lstlisting}[language=bash]
propwrap homework kerolox --name Ada
\end{lstlisting}
Produces summary markdown, assumptions text, CSV tables, and figures suitable
for a lab report without rebuilding a plotting stack from scratch.

% =============================================================================
\chapter{Limitations and responsible use}
\label{ch:limits}
% =============================================================================

\begin{enumerate}[leftmargin=*]
  \item \textbf{Ideal 1-D \cea{}} --- no kinetics, multiphase losses, injector
        efficiency, or nozzle divergence losses unless the user adds them.
  \item \textbf{Shifting/frozen brackets} --- flight sits elsewhere.
  \item \textbf{Efficiency knobs} --- rough, not a validated engine model.
  \item \textbf{Custom propellant cards} --- structure-checked, not
        thermodynamically certified by this project.
  \item \textbf{Licensing} --- MIT for \pw{}; review RocketCEA (GPL-family) and
        other dependencies before proprietary shipping.
  \item \textbf{Not flight-qualified software.}
  \item \textbf{No agency affiliation.}
\end{enumerate}

If a number will appear in a design review, label it \emph{theoretical CEA} or
\emph{estimated with $\eta$}, and keep hot-fire / flight data in a separate
column.

% =============================================================================
\chapter{Conclusions}
% =============================================================================

\pw{} packages a well-known physics engine (NASA~\cea{} via RocketCEA) behind a
modern Python interface that prefers SI units, typed results, and fair
propellant trades. The contribution is not a new chemistry model; it is
engineering hygiene: conversions that do not drift, identities that fail tests
when broken, and workflows that make correct comparisons convenient.

Version~0.1.x is suitable for education, early trade studies, and as a
propellant backbone under other tools. It is not a substitute for detailed
engine design analysis or certification evidence. With that boundary understood,
the library does the job it claims: give the same theoretical propellant numbers
tomorrow that it gave today, in units a SI codebase can trust.

% =============================================================================
\begin{thebibliography}{99}

\bibitem{gordon1994cea}
S.~Gordon and B.~J.~McBride,
\emph{Computer Program for Calculation of Complex Chemical Equilibrium
Compositions and Applications, I.~Analysis},
NASA Reference Publication 1311, 1994.
NTRS: \url{https://ntrs.nasa.gov/citations/19950013764}.

\bibitem{gordon1996cea}
S.~Gordon and B.~J.~McBride,
\emph{Computer Program for Calculation of Complex Chemical Equilibrium
Compositions and Applications, II.~Users Manual and Program Description},
NASA Reference Publication 1311, 1996.

\bibitem{rocketcea}
C.~Taylor / Applied Python,
\emph{RocketCEA} documentation and QuickStart goldens,
\url{https://rocketcea.readthedocs.io/}.

\bibitem{cantera}
D.~G.~Goodwin, H.~K.~Moffat, I.~Schoegl, R.~L.~Speth, and B.~W.~Weber,
\emph{Cantera: An object-oriented software toolkit for chemical kinetics,
thermodynamics, and transport processes},
\url{https://www.cantera.org}.

\bibitem{sutton2017}
G.~P.~Sutton and O.~Biblarz,
\emph{Rocket Propulsion Elements}, 9th~ed.,
John Wiley \& Sons, 2017.

\bibitem{huzel1992}
D.~K.~Huzel and D.~H.~Huang,
\emph{Modern Engineering for Design of Liquid-Propellant Rocket Engines},
AIAA Progress in Astronautics and Aeronautics, Vol.~147, 1992.

\bibitem{bipm_g0}
BIPM / CGPM,
conventional standard acceleration due to gravity
$\gzero=\SI{9.80665}{\metre\per\second\squared}$,
International System of Units (SI) documentation,
\url{https://www.bipm.org/}.

\bibitem{codata}
CODATA recommended values of the fundamental physical constants
(universal gas constant as used in software checks).

\bibitem{propwrap_repo}
S.~K.~V.,
\emph{propwrap} source repository,
\url{https://github.com/shanthoshkv/propwrap}.

\bibitem{propwrap_site}
S.~K.~V.,
\emph{propwrap} project website,
\url{https://aboutkvs.vercel.app/propwrap.html}.

\bibitem{propwrap_pypi}
\emph{propwrap} on the Python Package Index,
\url{https://pypi.org/project/propwrap/}.

\end{thebibliography}

% =============================================================================
\appendix
\chapter{Package layout}
\begin{verbatim}
propwrap/
  pyproject.toml
  README.md
  LICENSE
  src/propwrap/          # library
  tests/                 # pytest suite + validation_catalog.json
  docs/                  # user docs + this report
  examples/              # runnable scripts
  website/               # single-file HTML site
\end{verbatim}

\chapter{Document control}
\begin{center}
\begin{tabular}{@{}llp{7.5cm}@{}}
\toprule
Version & Date & Notes \\
\midrule
0.1 & 2026-07-30 & Initial technical report aligned with software v0.1.3 \\
\bottomrule
\end{tabular}
\end{center}

\vspace{1em}
\noindent\textit{End of report.}

\end{document}
